\documentclass[11pt]{article}
\usepackage[margin=1.1in]{geometry}
\usepackage{amsmath,amssymb,amsthm}
\usepackage{hyperref}

\theoremstyle{plain}
\newtheorem{theorem}{Theorem}[section]
\newtheorem{proposition}[theorem]{Proposition}
\newtheorem{lemma}[theorem]{Lemma}
\newtheorem{corollary}[theorem]{Corollary}
\newtheorem{conjecture}[theorem]{Conjecture}
\theoremstyle{definition}
\newtheorem{definition}[theorem]{Definition}
\newtheorem{remark}[theorem]{Remark}

\newcommand{\C}{\mathbb{C}}
\newcommand{\kk}{\Bbbk}
\newcommand{\Hess}{\operatorname{Hess}}
\newcommand{\Jac}{\operatorname{Jac}}
\newcommand{\adj}{\operatorname{adj}}
\newcommand{\JC}{\mathrm{JC}}
\newcommand{\HC}{\mathrm{HC}}

\title{The four-variable Hessian conjecture:\\ definitions, verified reductions, and the graded tower}
\author{HC$_4$ research project\thanks{All literature claims verified against
primary sources 2026-08-08/09 (see \texttt{status.md}); all computations in
exact rational arithmetic with fail-closed certificates (see
\texttt{certificates/}).}}
\date{August 9, 2026}

\begin{document}
\maketitle

\section{Definitions and precise problem statement}

Throughout, $\kk$ is a field of characteristic zero; the target problem is
over $\C$. For $f\in\kk[x_1,\dots,x_n]$ write $\Hess f=(\partial^2
f/\partial x_i\partial x_j)$ and $\nabla f=(\partial f/\partial
x_1,\dots,\partial f/\partial x_n)$. A polynomial map $F$ with $\Jac
F:=\det DF\in\kk^\times$ is a \emph{Keller map}.

\begin{definition}[formal Legendre transform; Meng 2006]
Let $\det\Hess f\in\kk^\times$. After an affine normalization, $p=\nabla
f(x)$ has a formal inverse $x(p)$ at the base point, and
$f^L(p):=\langle p,x(p)\rangle-f(x(p))$; equivalently $\nabla
f^L=(\nabla f)^{-1}$ as formal maps.
\end{definition}

\begin{conjecture}[$\HC_n$; Meng 2006]
If $f\in\kk[x_1,\dots,x_n]$ has $\det\Hess f\in\kk^\times$, then $f^L$ is a
polynomial.
\end{conjecture}

By the injectivity theorem below, over $\C$ this is equivalent to: every
\emph{gradient Keller map} $\nabla f$ is injective, equivalently a polynomial
automorphism. $\HC_4$ is the target problem.

\begin{remark}[a different ``Hessian conjecture'']
de Bondt--van den Essen (Proc.\ AMS 133 (2005) 2201--2205) use ``HC$(n)$''
for the \emph{nilpotent} statement ($\Hess f$ nilpotent $\Rightarrow$
$x+\nabla f$ invertible). Their Theorem~1.1 proves nilpotent-HC$(2n)$
$\Rightarrow$ $\JC_n$ for nilpotent Keller maps, and the nilpotent
\emph{symmetric} case is settled for $n\le4$ (any $H$) and $n=5$
homogeneous (JPAA 196 (2005) 135--148). None of this settles Meng's
$\HC_4$.
\end{remark}

\subsection{Classical inputs (all source-verified)}

\begin{theorem}[injectivity $\Rightarrow$ automorphism; BBR 1962, Ax 1969,
Cynk--Rusek 1991, BCW Thm 2.1]\label{thm:inj}
An injective polynomial map $\C^n\to\C^n$ is bijective and its inverse is
polynomial.
\end{theorem}

\begin{theorem}[Wang 1980; proof of Oda, cf.\ BCW Thm 2.4]\label{thm:wang}
Every Keller map of degree $\le2$ (any $n$, char $\ne2$, not necessarily
homogeneous) is invertible.
\end{theorem}

\begin{theorem}[Dillen 1991 ($n\le2$); de Bondt, JPAA 219 (2015), Thm 2.1
($n\le3$)]\label{thm:hc3}
For $\kk$ of characteristic zero and $f\in\kk[x_1,\dots,x_n]$ with $n\le3$:
$\nabla f$ satisfies the Jacobian conjecture; in the Keller case the inverse
is polynomial. Moreover if $3\le d=\deg f$ and $\det\Hess f\in\kk$, some
$T\in GL_n(\kk)$ makes $\Hess(f(Tx))$ vanish below the anti-diagonal.
Both structural tools (anti-triangularization and the weight theorem) fail
for $n\ge4$ by explicit counterexamples in the same paper.
\end{theorem}

\begin{theorem}[Gordan--Noether 1876; Watanabe--de Bondt, arXiv:1703.07624]
\label{thm:gn}
A form in $n\le4$ variables (alg.\ closed, char $0$) with identically
vanishing Hessian determinant becomes, after a linear change, a form in
$\le n-1$ variables. For $n=5$ this fails (Perazzo cubic
$x_1^2x_3+x_1x_2x_4+x_2^2x_5$); the $n=5$ vanishing-Hessian forms properly
involving 5 variables are, after a linear change, the homogeneous elements
of $K[x_1,x_2][\Delta]$, $\Delta=p_3(x_1,x_2)x_3+p_4(x_1,x_2)x_4+p_5(x_1,x_2)x_5$.
\end{theorem}

\section{The bridges and stabilization (proofs re-verified)}

\begin{proposition}[Meng 2006]\label{prop:bridges}
(i) $\JC_n\Rightarrow\HC_n$.
(ii) $\HC_{2n}\Rightarrow\JC_n$: for a Keller map $F$, the doubling
$\varphi(x,y)=\langle y,F(x)\rangle$ has
$\det\Hess\varphi=(-1)^n(\Jac F)^2$ and $\nabla\varphi$ triangular, so a
polynomial $\varphi^L$ yields a polynomial $F^{-1}$.
(iii) If $\det\Hess f\in\kk^\times$ and $\nabla f$ is not injective, $f^L$
is not a polynomial. In particular $\boxed{\HC_4\Rightarrow\JC_2}$.
\end{proposition}

\begin{lemma}[stabilization]\label{lem:stab}
$\JC_n$ false $\Rightarrow\JC_{n+1}$ false (pad by identity); $\HC_n$ false
$\Rightarrow\HC_{n+1}$ false ($f\oplus\tfrac12t^2$; the Legendre transform
separates).
\end{lemma}

\section{Status inputs (independently machine-verified here)}

Alp\"oge's $F$ ($\Jac\equiv-2$, three-point fiber) refutes $\JC_3$;
Meng--Yang's $\Psi=A^2+13A+2B$ ($\deg 14$, $\det\Hess\Psi\equiv128$,
gradient collision at $(\pm1,\mp\tfrac32,0,0,0)$) refutes $\HC_5$; the full
Schur family satisfies $\det\Hess(B+\tfrac\lambda2A^2+\mu A)\equiv4\lambda$.
Certificates: \texttt{verify\_alpoge\_3d.py},
\texttt{verify\_meng\_yang\_hc5.py}, \texttt{verify\_mengyang\_fast.py}.
With Lemma~\ref{lem:stab} and Theorem~\ref{thm:hc3} this leaves exactly
$\JC_2$ and $\HC_4$ open.

\section{The graded tower for $\det\Hess f\in\C^\times$}

Let $f=\sum_k f_k$ (homogeneous parts), $d=\deg f$, $n=4$ unless stated.

\begin{lemma}[T0]\label{lem:t0}
$[\det\Hess f]_0=\det\Hess f_2$. Hence the quadratic part is nondegenerate:
$\det\Hess f_2=c$.
\end{lemma}

\begin{lemma}[T1]\label{lem:t1}
If $d\ge3$ then $\det\Hess f_d\equiv0$ (the top graded piece, degree
$n(d-2)>0$, of a constant). With Theorem~\ref{thm:gn}: \emph{the leading
form of any $\HC_4$ candidate is, WLOG, in $\C[x_1,x_2,x_3]$.}
\end{lemma}

\begin{lemma}[T3]\label{lem:t3}
If $f_d\in\C[x_1,x_2,x_3]$, $d\ge3$, then
$[\det\Hess f]_{4d-9}=\det{}_3(\Hess_3f_d)\cdot\partial^2f_{d-1}/\partial
x_4^2=0$: either $f_d$ degenerates to $\C[x_1,x_2]$, or $f_{d-1}$ has
$x_4$-degree $\le1$.
\end{lemma}

\begin{proof}[Proof sketch for T0/T1/T3]
Each entry of $\Hess f_k$ is homogeneous of degree $k-2$; a permutation
product mixes graded pieces additively. For T3: entries of $\Hess f_d$
vanish in row/column 4, so a product of degree $3(d-2)+(d-3)$ must place
its unique lower-degree entry at position $(4,4)$; summing over
$3$-permutations of the upper block gives the $3\times3$ determinant.
Machine spot-checks: \texttt{src/hc4lib.py} self-tests.
\end{proof}

\begin{lemma}[T4; midpoint identity, $\deg f\le4$; fully symbolic machine
proof]\label{lem:t4}
For $\deg f\le4$ and all $m,v$:
$\nabla f(m+v)-\nabla f(m-v)=2[\Hess f(m)\,v+\nabla f_4(v)]$. Hence $\nabla
f$ is non-injective iff $\exists\,v\ne0,m$ with $\Hess f(m)v=-\nabla
f_4(v)$; and for $\deg f\le3$, $\HC_n$ holds for all $n$ (Wang's case:
$\Hess f(m)v=0$ forces $v=0$).
\end{lemma}

\section{The quadratic-pivot theorem and the affine-pivot reduction}

\begin{theorem}[Q2; proved here]\label{thm:q2}
Let $f=e_0(x')+x_4e_1(x')+\tfrac\sigma2x_4^2$ with $x'=(x_1,x_2,x_3)$,
$\sigma\in\C^\times$, $\deg e_1\le2$, and $\det\Hess f=c\in\C^\times$. Set
$\tilde e_0:=e_0-e_1^2/(2\sigma)$ and $K:=\Hess_3e_1$ (constant). Then
\[
\det\Hess f=\sigma\,\det{}_3\!\bigl(\Hess_3\tilde e_0+(x_4+e_1(x')/\sigma)K\bigr),
\]
and since $u=x_4+e_1(x')/\sigma$ sweeps $\C$ for fixed $x'$, the pencil
identity $\det_3(\Hess_3\tilde e_0+sK)\equiv c/\sigma$ holds for all
$s\in\C$. Any collision $\nabla f(p)=\nabla f(q)$ forces a common
$\lambda=p_4+e_1(p')/\sigma=q_4+e_1(q')/\sigma$ and a collision of
$\nabla'(\tilde e_0+\lambda e_1)$; by the pencil identity and
Theorem~\ref{thm:hc3} this map is injective, so $p=q$. \emph{Hence $\HC_4$
holds for every polynomial that, in some affine coordinates, is quadratic
in a variable with constant leading coefficient and $\deg e_1\le2$.}
\end{theorem}

\begin{theorem}[T5; affine-pivot reduction in degree 4 --- proved modulo the
branch closures listed in \texttt{lemma\_ledger.md}]\label{thm:t5}
Let $\deg f=4$, $\det\Hess f\in\C^\times$. Then, after an affine change,
either $f$ satisfies the hypotheses of Theorem~\ref{thm:q2} (and $\HC_4$
holds for it), or
\[
f=e_0(x_1,x_2,x_3)+x_4\,e_1(x_1,x_2,x_3)\qquad(\textbf{class } AP_4),
\]
with $\deg e_0\le4$, $\deg e_1\le2$, and, writing $g=\nabla'e_1$,
$P=\Hess_3e_0$, $K=\Hess_3e_1$,
\[
\text{(c0) } g^\top(\adj P)g\equiv-c,\qquad
\text{(c1) } g^\top L(P,K)g\equiv0,\qquad
\text{(c2) } g^\top(\adj K)g\equiv0,
\]
where $\adj(P+tK)=\adj P+tL(P,K)+t^2\adj K$.
\end{theorem}

\begin{remark}
The class $AP_4$ contains all Meng doublings
$\langle(x_3,x_4),F(x_1,x_2)\rangle$ of planar Keller maps (e.g.\ de
Bondt's anti-triangularization counterexample), whose injectivity is
equivalent to $\JC_2$ instances. Thus the affine-pivot class is exactly
where $\JC_2$ embeds into $\HC_4$, and no argument that closes all of
$AP_4$ can avoid proving planar-Jacobian statements. This calibrates the
difficulty: full $\HC_4$ likely requires either $\JC_2$ or a structure
theory of $AP_4$ modulo doublings.
\end{remark}

\section{Obstruction results for the known counterexample mechanisms}

\begin{theorem}[D1/D1$^+$; proved here, machine certificate]\label{thm:d1}
For every $\lambda\ne0,\mu$, the Meng--Yang potential
$\psi_{\lambda,\mu}=B+\tfrac\lambda2A^2+\mu A$ admits no $v\in\C^5\setminus
\{0\}$ with $D_v^2\psi\equiv\gamma$ for any constant $\gamma$ (including
$0$). Hence no affine coordinates exhibit an affine pivot (Schur descent)
or a quadratic pivot on any member of the known $\HC_5$ counterexample
family: \emph{the known descent mechanism cannot be iterated to reach
$\HC_4$.} (Certificate: 95 homogeneous $v$-quadratics from the
degree-$\ge6$ coefficients; Gr\"obner leading terms contain a pure power of
each $v_i$.)
\end{theorem}

\begin{theorem}[G1; proved here, machine certificate]\label{thm:g1}
For $F\in\{$Alp\"oge $F$, Gao $F_4$, Gao $F_5\}$ the space $\{W:W\cdot
J_F(x)\text{ symmetric }\forall x\}$ is $\{0\}$. Since $L_1\circ F\circ
L_2$ is a gradient map for affine $L_i$ iff some invertible $W$
symmetrizes $J_F$ (take $W=(T^\top)^{-1}S$; conversely the polynomial
Poincar\'e lemma), \emph{neither known dimension-4 Jacobian counterexample
is affinely equivalent to a gradient map}. The Alp\"oge case is the
consistency control demanded by Theorem~\ref{thm:hc3}.
\end{theorem}

\section{Dependency graph}

\begin{verbatim}
                    [BBR/Ax/Cynk-Rusek]   [Wang/Oda]     [Gordan-Noether]
                          |                   |                |
                        (L-INJ)             (T4)             (T2)
                          |                   |                |
 [de Bondt HC_3] ---------+---------+       (deg<=3 case)    (T1,T3 tower)
        |                           |                          |
      (Q2) <---- pencil identity ---+                        (T5 reduction)
        |                                                      |
        +------------------ HC_4(deg 4) = Q2-side + AP_4 -----+
                                              |
 [Meng doubling] --- AP_4 contains doubled planar Keller maps --- [JC_2]
 [Meng-Yang Schur descent] --(D1/D1+: no pivot on known family)--X--> HC_4
 [Gao F_4,F_5] -----(G1: not affinely gradient-equivalent)-----X--> HC_4
\end{verbatim}

Statements in square brackets are literature (source-verified); parenthesized
statements are proved in this project (certificates in
\texttt{certificates/}); X marks proved obstructions.

\end{document}
